\documentclass{article}
\pdfinfoomitdate=1
\pdftrailerid{}

\usepackage{float}
\restylefloat{table}

\usepackage{booktabs}

\title{Team Contributions: Final\\\progname}

\author{\authname}

\date{March 25, 2026}

\input{../Comments}
\input{../Common}

\begin{document}

\maketitle

This document summarizes the contributions of each team member for the final
demonstration and documentation. The time period of interest is the time
between Rev 0 and the Final documentation; the contributions prior to Rev0 are
NOT included.

\section{Team Meeting Attendance}

\wss{For each team member how many team meetings have they attended over the
  time period of interest.  This number should be determined from the meeting
  issues in the team's repo.  The first entry in the table should be the total
  number of team meetings held by the team.}

\begin{table}[H]
  \centering
  \begin{tabular}{ll}
    \toprule
    \textbf{Student} & \textbf{Meetings} \\
    \midrule
    Total            & 11                \\
    Xiaotian Lou     & 11                \\
    Zifan Si         & 11                \\
    Jianqing Liu     & 11                \\
    Shike Chen       & 11                \\
    \bottomrule
  \end{tabular}
\end{table}

\wss{If needed, an explanation for the counts can be provided here.}

\section{Supervisor/Stakeholder Meeting Attendance}

\wss{For each team member how many supervisor/stakeholder team meetings have
  they attended over the time period of interest.  This number should be determined
  from the supervisor meeting issues in the team's repo.  The first entry in the
  table should be the total number of supervisor and team meetings held by the
  team.  If there is no supervisor, there will usually be meetings with
  stakeholders (potential users) that can serve a similar purpose.}

\noindent \textbf{Supervisor's Name: } Sirouspour, Shahin

\begin{table}[H]
  \centering
  \begin{tabular}{ll}
    \toprule
    \textbf{Student} & \textbf{Meetings} \\
    \midrule
    Total            & 0                 \\
    Xiaotian Lou     & 0                 \\
    Zifan Si         & 0                 \\
    Jianqing Liu     & 0                 \\
    Shike Chen       & 0                 \\
    \bottomrule
  \end{tabular}
\end{table}

\wss{If needed, an explanation for the counts can be provided here.}

\section{Lecture Attendance}

\wss{For each team member how many lectures have they attended over the time
  period of interest.  This number should be determined from the lecture issues in
  the team's repo. You can find the number of lectures in the time period of
  interest by looking at the
  \href{https://calendar.google.com/calendar/u/0/embed?src=rnboqiaki1k2la7rpu3bn0um58@group.calendar.google.com&ctz=America/Toronto}
  {Google calendar} for the capstone course.}

\wss{NOTE: There will be approximately 1 lecture between the Rev0 and Rev1
  demos}

\begin{table}[H]
  \centering
  \begin{tabular}{ll}
    \toprule
    \textbf{Student} & \textbf{Lectures} \\
    \midrule
    Total            & 1                 \\
    Xiaotian Lou     & 0                 \\
    Zifan Si         & 0                 \\
    Jianqing Liu     & 0                 \\
    Shike Chen       & 0                 \\
    \bottomrule
  \end{tabular}
\end{table}

\wss{If needed, an explanation for the lecture attendance can be provided here.}

\section{TA Document Discussion Attendance}

\wss{For each team member how many of the informal document discussion meetings
  with the TA were attended over the time period of interest.}

\noindent \textbf{TA's Name: } Tanya Djavaherpour

\begin{table}[H]
  \centering
  \begin{tabular}{ll}
    \toprule
    \textbf{Student} & \textbf{Lectures} \\
    \midrule
    Total            & 2                 \\
    Xiaotian Lou     & 2                 \\
    Zifan Si         & 2                 \\
    Jianqing Liu     & 2                 \\
    Shike Chen       & 2                 \\
    \bottomrule
  \end{tabular}
\end{table}

\wss{If needed, an explanation for the attendance can be provided here.}

\section{Commits}

\wss{For each team member how many commits to the main branch have been made
  over the time period of interest.  The total is the total number of commits for
  the entire team since the beginning of the term.  The percentage is the
  percentage of the total commits made by each team member.}

\begin{table}[H]
  \centering
  \begin{tabular}{lll}
    \toprule
    \textbf{Student} & \textbf{Commits} & \textbf{Percent} \\
    \midrule
    Total            & 368              & 100.0\%          \\
    Xiaotian Lou     & 97               & 26.4\%           \\
    Jianqing Liu     & 96               & 26.1\%           \\
    Zifan Si         & 108              & 29.3\%           \\
    Shike Chen       & 67               & 18.2\%           \\
    \bottomrule
  \end{tabular}
\end{table}

\noindent\textit{Notes.} Commit counts reflect merges to the \texttt{main} branch as well as contributions in experimental branches (including Copilot-assisted branches such as \texttt{copilot/sub-pr-302}, \texttt{copilot/sub-pr-302-again}, and \texttt{copilot/explore-repo-and-pr-solution}) attributed to the team member who initiated them. During the Final period, Shike Chen focused primarily on the labelling application, SAM backend integration, and data pipeline tooling, while Xiaotian Lou concentrated on unit testing, ML model benchmarking, and TensorRT inference optimization. Much of this work involved on-device iterations, training runs, and tooling validation that lived in feature branches or local workspaces before being upstreamed. As a result, the raw \texttt{main}-branch totals alone would under-represent their effort; the adjusted figures above provide a more balanced picture.

\section{Issue Tracker}

\wss{For each team member how many issues have they authored (including open and
  closed issues (O+C)) and how many have they been assigned (only counting closed
  issues (C only)) over the time period of interest.}

\begin{table}[H]
  \centering
  \begin{tabular}{lll}
    \toprule
    \textbf{Student} & \textbf{Authored (O+C)} & \textbf{Assigned (C only)} \\
    \midrule
    Xiaotian Lou     & 4                       & 3                          \\
    Zifan Si         & 4                       & 5                          \\
    Jianqing Liu     & 31                      & 11                         \\
    Shike Chen       & 9                       & 8                          \\
    \bottomrule
  \end{tabular}
\end{table}

\noindent\textit{Notes.}
\begin{itemize}
  \item \textbf{Counting rule:} "Authored (O+C)" counts issues opened by a member (open or closed). "Assigned (C only)" counts issues where the member was the assignee at closure.
  \item \textbf{Role-based workflow:} During the Final period, Jianqing Liu (project manager, "Pega") centrally opened, assigned, and closed most issues based on team decisions captured in meeting notes; hence his "Authored" count is higher.
  \item \textbf{Credit vs.\ logistics:} Centralized opening/closing is administrative. Actual implementation credit is reflected in commits, PRs, and code reviews; an issue can have multiple contributors beyond the final assignee.
  \item \textbf{Assignment semantics:} Ownership may change during an issue's lifetime; we attribute "Assigned (C only)" to the final assignee at closure to represent delivery responsibility.
\end{itemize}

\section{CICD}

\wss{Say how CICD is used in your project}

We use GitHub Actions with five workflows that cover documentation, firmware,
backend, and the operator UI.

\paragraph{build-tex (push to \texttt{main}, \texttt{docs/**}, or manual)}
Compiles LaTeX with TeX~Live~2025 on Ubuntu and commits the built PDFs back to the repository.
The job runs \texttt{make -B} under \texttt{docs/} and then auto-commits any updated \texttt{.pdf} files
(\texttt{contents: write}).

\paragraph{check-tex (pull requests touching \texttt{docs/**})}
Validates that the LaTeX compiles on PRs, uploads any changed PDFs as an artifact, then formats all
\texttt{.tex} sources with \texttt{latexindent} (80-column wrapping) and commits the formatting changes.
This keeps the documentation reproducible and consistently formatted.

\paragraph{firmware-ci (push/PR under \texttt{src/pcb-firmware/**})}
Builds the embedded firmware with Cargo on Ubuntu and uploads the resulting binary as an artifact.
The job also enforces Rust source formatting via \texttt{cargo fmt --check}.

\paragraph{backend-ci (pull requests touching \texttt{src/backend/**})}
Runs backend lint checks on Ubuntu using Python~3.10 with pip dependency caching. The workflow installs dependencies from \texttt{src/backend/requirement.txt} (excluding any pinned \texttt{pip} line) and performs a basic syntax check via \texttt{python -m compileall -q .} in \texttt{src/backend}. This helps catch Python syntax errors early before merging.

\paragraph{react-ci (push/PR for the React app)}
Builds and smoke-tests the operator UI on Node~20. Dependencies are installed
with \texttt{npm ci}; a lightweight \texttt{ci:smoke} runs on every change.
Unit tests and production build currently run in \emph{best-effort} mode
(non-blocking), and any build artifacts are uploaded when present.

\section{Team Charter Trigger Items}

\subsection*{Quantified Triggers (per Team Charter)}
\begin{itemize}
  \item \textbf{Meeting cadence \& attendance:} weekly team meetings; prior notice required for any absence; target attendance $\geq$85\%.
  \item \textbf{Preparation \& quality:} come to meetings with concise updates, blockers, and options; maintain decision logs and lightweight RACI notes.
  \item \textbf{PR/issue hygiene:} every task has a GitHub issue with an owner and estimate; all code changes go through PRs that reference an issue; no direct pushes to \texttt{main}.
  \item \textbf{Review SLA \& CI gates:} at least one reviewer approval; address review comments within 7 days (unless a different SLA is agreed); CI must be green (docs build, firmware build, UI smoke).
  \item \textbf{Documentation standards:} LaTeX compiles reproducibly on CI; PDFs are published on \texttt{main}; \texttt{latexindent} enforces consistent formatting on PRs.
  \item \textbf{Runtime KPIs (demo readiness):} track end-to-end latency (p50/p95), FPS, target re-acquisition time, crash-free session rate, and UI task success; for the final demo we target $\sim$100--500\,ms latency, $\geq$15\,FPS, and $\leq$2\,s re-acquisition.
  \item \textbf{Underperformance trigger:} if a member underdelivers for $>$3 weeks, initiate pairing/guardrails and, if needed, escalate to TA/instructor.
\end{itemize}

\subsection*{Violations Observed During the Final Period}
None observed. All members met the charter triggers: weekly meetings were held
with advance notice for any conflicts; PRs referenced issues and passed CI;
review comments were addressed within the 7-day SLA; documentation built
reproducibly; runtime KPIs were recorded for the demo.

\subsection*{Plan to Sustain Compliance}
\begin{itemize}
  \item Keep branch protection (no direct pushes to \texttt{main}; $\geq$1 review; CI
        required).
  \item Maintain formatting gates: \texttt{latexindent} for LaTeX, \texttt{cargo fmt}
        for Rust firmware, and lint checks for Python and React sources.
  \item Continue using meeting-issue templates (owner, ETA, blockers) and maintain the
        decision log to preserve cadence and traceability.
  \item Monitor runtime KPI thresholds (p50/p95 latency, FPS, re-acquisition time) on
        each demo and EXPO run.
\end{itemize}

\section{Additional Productivity Metrics}

\wss{If your team has additional metrics of productivity, please feel free to
  add them to this report.}
At this time, we have no additional productivity metrics to report.
\end{document}
